\documentclass[11pt]{article}

\usepackage[T1]{fontenc}
\usepackage{amsmath,amssymb,amsthm}
\usepackage[margin=1in]{geometry}
\usepackage[colorlinks=true,linkcolor=blue,urlcolor=blue]{hyperref}

\newtheorem{theorem}{Theorem}
\newtheorem{lemma}[theorem]{Lemma}
\newtheorem{corollary}[theorem]{Corollary}
\newtheorem{proposition}[theorem]{Proposition}
\newtheorem{definition}[theorem]{Definition}
\newtheorem{conjecture}[theorem]{Conjecture}
\theoremstyle{remark}
\newtheorem{remark}[theorem]{Remark}

\newcommand{\R}{\mathbb R}
\newcommand{\norm}[1]{\left\|#1\right\|}
\newcommand{\ip}[2]{\left\langle#1,#2\right\rangle}
\newcommand{\cK}{\mathcal K}
\newcommand{\esssup}{\operatorname*{ess\,sup}}
\setlength{\emergencystretch}{2em}

\title{The q4 Cross Defect Forces a Moving Frequency Corridor}
\author{John Lawson and Codex}
\date{25 July 2026}

\begin{document}
\maketitle

\begin{abstract}
This round tests whether the pressure-free primal--adjoint commutator
anomaly, used through available Lorentz norm bounds, forces the summable
amplitude anti-correlation sought in the preceding round.  It does not.
Instead, the terminal cross defect defines a moving dyadic frequency at
which half of its conserved pairing first appears.  The weak-\(L^3\)
clock limits the speed of that frequency, while the unresolved
high-frequency pairing forces a product of primal and adjoint enstrophy.

For the exact q4 clock, the resulting pointwise lower bound is
\[
 \norm{\nabla u(T-\tau)}_2
 \norm{\nabla a(T-\tau)}_2
 \gtrsim
 \tau^{-9/11}
 \log(e+\tau_0/\tau)^{2/11}.
\]
The product of the two terminal dissipation tails is therefore at least
of order
\(\tau^{4/11}\log(e+\tau_0/\tau)^{4/11}\).

A refined commutator estimate gives an exact alternative: the cross
frequency either reaches the q4 nonlinear clock or remains in a
sub-clock corridor paid by accumulated \(M\norm{\nabla a}_2\).
An explicit scalar temporal ledger saturates the corridor, every energy
budget, and the divergent weighted adjoint dissipation.  It is not a
Navier--Stokes construction.  The live task is now to exclude corridor
reuse using genuinely same-trajectory phase, sign, residence, or
projected-pressure structure.  No Clay alternative is proved.
\end{abstract}

\section{Epistemic status and exact input}

All theorems below are conditional on the repository's energy-efficient
exact q4 branch and projected-Oseen entrance.  Those hypotheses are not
derived for arbitrary Clay data, and external review is pending.

Let \(T\) be the putative first singular time.  The input supplies smooth
preterminal \(u\), an energy-class projected adjoint \(a\), and constants
\[
 U:=\sup_{t<T}\norm{u(t)}_2<\infty,
 \qquad
 A:=\sup_{t<T}\norm{a(t)}_2<\infty,
\]
such that
\begin{equation}
\label{eq:pairing}
 \ip{u(t)}{a(t)}=c_0>0
 \qquad(0<t<T),
\end{equation}
\[
 u,a\in L^2_t\dot H^1_x,
 \qquad
 a(t)\rightharpoonup0
 \quad(t\uparrow T).
\]

For reverse terminal time \(0<\tau<\tau_0\), put
\begin{equation}
\label{eq:MXY}
 M(\tau):=\norm{u(T-\tau)}_{L^{3,\infty}},
 \quad
 X(\tau):=\norm{\nabla u(T-\tau)}_2,
 \quad
 Y(\tau):=\norm{\nabla a(T-\tau)}_2.
\end{equation}
The exact clock gives, for large \(\lambda\),
\begin{equation}
\label{eq:clock}
 \left|
 \left\{\tau:M(\tau)>\lambda\right\}
 \right|
 \lesssim
 \frac{\lambda^{-11/2}}{\log(e+\lambda)}.
\end{equation}

Choose a real radial
\(\psi\in C_c^\infty(\R^3)\) satisfying
\[
 \psi(\xi)=1\quad(|\xi|\leq1),
 \qquad
 \psi(\xi)=0\quad(|\xi|\geq2),
\]
and set
\[
 B_N:=\psi(D/N).
\]
The pressure-free identity from the preceding round is
\begin{equation}
\label{eq:comm-identity}
 \frac d{dt}\ip{u}{B_Na}
 =
 \ip{[(u\cdot\nabla),B_N]u}{a}.
\end{equation}
For each fixed \(N\), its terminal pairing vanishes.  Define
\begin{equation}
\label{eq:KN}
 K_N(\tau)
 :=
 \ip{u(T-\tau)}{B_Na(T-\tau)}.
\end{equation}
Then
\begin{equation}
\label{eq:KN-limits}
 \lim_{\tau\downarrow0}K_N(\tau)=0,
 \qquad
 \lim_{N\to\infty}K_N(\tau)=c_0
 \quad(\tau>0).
\end{equation}

The Clay problem remains unsolved.

\section{The moving cross-pairing frequency}

\begin{definition}
\label{def:kappa}
For almost every \(\tau>0\), define
\begin{equation}
\label{eq:kappa}
 \boxed{
 \kappa(\tau)
 :=
 \min
 \left\{
 2^k:K_{2^k}(\tau)\geq c_0/2
 \right\}.
 }
\end{equation}
\end{definition}

The two limits in \eqref{eq:KN-limits} show that this dyadic minimum is
well-defined.

Write
\begin{equation}
\label{eq:B1B2H}
 B_1(\tau):=\int_0^\tau M(s)\,ds,
 \quad
 B_2(\tau):=\int_0^\tau M(s)^2\,ds,
 \quad
 H(\tau):=\int_0^\tau M(s)Y(s)\,ds.
\end{equation}
All fields under these integrals are evaluated at \(T-s\).

\begin{lemma}[Coarse pressure-free speed limit]
\label{lem:coarse-speed}
For every \(N>0\),
\begin{equation}
\label{eq:coarse-speed}
 \left|
 \ip{[(u\cdot\nabla),B_N]u}{a}
 \right|
 \lesssim
 UA\,N^2M.
\end{equation}
\end{lemma}

\begin{proof}
Put \(T_u:=u\cdot\nabla\).  Lorentz H\"older and band-limited
Bernstein give
\[
 \begin{aligned}
 |\ip{T_uB_Nu}{a}|
 &\leq
 \norm{u}_{L^{3,\infty}}
 \norm{\nabla B_Nu}_{L^{6,2}}
 \norm{a}_2\\
 &\lesssim
 M N^2UA.
 \end{aligned}
\]
For the second commutator term, self-adjointness and
\(\nabla\cdot u=0\) give
\[
 \begin{aligned}
 |\ip{B_NT_uu}{a}|
 &=
 |\ip{T_uu}{B_Na}|\\
 &=
 \left|
 \int u_i u_j\partial_jB_Na_i
 \right|\\
 &\lesssim
 \norm{u}_{L^{3,\infty}}\norm{u}_2
 \norm{\nabla B_Na}_{L^{6,2}}\\
 &\lesssim
 M N^2UA.
 \end{aligned}
\]
The viscous terms and both pressures have already cancelled in
\eqref{eq:comm-identity}.
\end{proof}

\begin{theorem}[Cross-frequency motion]
\label{thm:cross-motion}
For almost every sufficiently small \(\tau\),
\begin{equation}
\label{eq:kappa-B1}
 \boxed{
 \kappa(\tau)^2B_1(\tau)\gtrsim1,
 }
\end{equation}
\begin{equation}
\label{eq:XY-kappa}
 \boxed{
 X(\tau)Y(\tau)\gtrsim\kappa(\tau)^2,
 }
\end{equation}
and hence
\begin{equation}
\label{eq:XY-B1}
 \boxed{
 X(\tau)Y(\tau)\gtrsim B_1(\tau)^{-1}.
 }
\end{equation}
\end{theorem}

\begin{proof}
Integrating \eqref{eq:comm-identity} to the terminal boundary and using
\eqref{eq:KN-limits} gives
\[
 K_N(\tau)
 =
 -\int_{T-\tau}^T
 \ip{[(u\cdot\nabla),B_N]u}{a}\,dt.
\]
Lemma~\ref{lem:coarse-speed} therefore yields
\[
 |K_N(\tau)|
 \lesssim
 UA\,N^2B_1(\tau).
\]
At \(N=\kappa(\tau)\), the left side is at least \(c_0/2\).
This proves \eqref{eq:kappa-B1}.

Minimality in Definition~\ref{def:kappa} gives
\[
 K_{\kappa(\tau)/2}(\tau)<c_0/2.
\]
Subtract this from the full pairing \eqref{eq:pairing}:
\begin{equation}
\label{eq:high-cross}
 \frac{c_0}{2}
 <
 \ip{u(T-\tau)}
 {(I-B_{\kappa(\tau)/2})a(T-\tau)}.
\end{equation}
The multiplier in \eqref{eq:high-cross} vanishes on a ball of radius
comparable to \(\kappa\).  Plancherel gives
\[
 \left|
 \ip{u}{(I-B_{\kappa/2})a}
 \right|
 \lesssim
 \kappa^{-2}\norm{\nabla u}_2\norm{\nabla a}_2
 =
 \kappa^{-2}XY.
\]
This proves \eqref{eq:XY-kappa}; combining it with
\eqref{eq:kappa-B1} proves \eqref{eq:XY-B1}.
\end{proof}

\section{The logarithmic q4 consequence}

Set
\[
 \ell(\tau):=\log(e+\tau_0/\tau).
\]

\begin{lemma}[Integrated q4 clock]
\label{lem:clock-integrals}
For all sufficiently small \(\tau\),
\begin{equation}
\label{eq:B1B2-clock}
 \boxed{
 B_1(\tau)
 \lesssim
 \tau^{9/11}\ell(\tau)^{-2/11},
 \qquad
 B_2(\tau)
 \lesssim
 \tau^{7/11}\ell(\tau)^{-4/11}.
 }
\end{equation}
\end{lemma}

\begin{proof}
The decreasing rearrangement consequence of \eqref{eq:clock} is
\[
 M^*(s)
 \lesssim
 1+s^{-2/11}\ell(s)^{-2/11}.
\]
Hardy--Littlewood and integration over \(0<s<\tau\) give the first
estimate in \eqref{eq:B1B2-clock}.  Squaring the rearrangement gives the
second.  A dyadic decomposition of \((0,\tau)\) controls the slowly
varying logarithm.
\end{proof}

\begin{corollary}[Pointwise joint-enstrophy floor]
\label{cor:pointwise-floor}
For almost every sufficiently small \(\tau\),
\begin{equation}
\label{eq:pointwise-floor}
 \boxed{
 X(\tau)Y(\tau)
 \gtrsim
 \tau^{-9/11}\ell(\tau)^{2/11}.
 }
\end{equation}
\end{corollary}

\begin{proof}
Combine Theorem~\ref{thm:cross-motion} and
Lemma~\ref{lem:clock-integrals}.
\end{proof}

Define the two terminal dissipation tails
\begin{equation}
\label{eq:PD}
 P(\tau):=\int_0^\tau X(s)^2\,ds,
 \qquad
 D(\tau):=\int_0^\tau Y(s)^2\,ds.
\end{equation}

\begin{corollary}[Joint terminal tail floor]
\label{cor:tail-floor}
For all sufficiently small \(\tau\),
\begin{equation}
\label{eq:PD-floor}
 \boxed{
 P(\tau)D(\tau)
 \gtrsim
 \tau^{4/11}\ell(\tau)^{4/11}.
 }
\end{equation}
Consequently each of \(P,D\) has a lower floor of the displayed order,
with a constant depending on the finite total tail of the other field.
\end{corollary}

\begin{proof}
Integrating \eqref{eq:pointwise-floor} over
\([\tau/2,\tau]\) gives
\[
 \int_0^\tau XY
 \gtrsim
 \tau^{2/11}\ell(\tau)^{2/11}.
\]
Cauchy--Schwarz gives
\[
 \int_0^\tau XY
 \leq
 P(\tau)^{1/2}D(\tau)^{1/2}.
\]
Squaring proves \eqref{eq:PD-floor}.  The individual statements follow
from the finite upper bounds \(P(\tau_0),D(\tau_0)<\infty\).
\end{proof}

\begin{corollary}[Inverse-clock condition]
\label{cor:inverse-clock}
Every entrance satisfies
\begin{equation}
\label{eq:inverse-clock}
 \boxed{
 \int_0^{\tau_0}\frac{d\tau}{B_1(\tau)}<\infty.
 }
\end{equation}
\end{corollary}

\begin{proof}
Theorem~\ref{thm:cross-motion} and Cauchy--Schwarz imply
\[
 \int_0^{\tau_0}\frac{d\tau}{B_1(\tau)}
 \lesssim
 \int_0^{\tau_0}XY\,d\tau
 \leq
 P(\tau_0)^{1/2}D(\tau_0)^{1/2}.
\]
\end{proof}

\section{Clock branch versus cross-frequency corridor}

\begin{lemma}[Refined pressure-free speed limit]
\label{lem:refined-speed}
For every \(N>0\),
\begin{equation}
\label{eq:refined-speed}
 \boxed{
 \left|
 \ip{[(u\cdot\nabla),B_N]u}{a}
 \right|
 \lesssim
 A N^{3/2}M^2
 +
 U NMY.
 }
\end{equation}
\end{lemma}

\begin{proof}
Skew-adjointness of \(T_u\) gives
\[
 \begin{aligned}
 |\ip{T_uB_Nu}{a}|
 &=
 |\ip{B_Nu}{T_ua}|\\
 &\lesssim
 \norm{B_Nu}_{L^{6,2}}
 \norm{u}_{L^{3,\infty}}
 \norm{\nabla a}_2\\
 &\lesssim
 UNMY.
 \end{aligned}
\]
For the other term,
\[
 |\ip{B_NT_uu}{a}|
 =
 \left|
 \int u_i u_j\partial_jB_Na_i
 \right|.
\]
O'Neil H\"older gives
\[
 \norm{u\otimes u}_{L^{3/2,\infty}}\lesssim M^2.
\]
The band-limited estimate
\[
 \norm{\nabla B_Na}_{L^{3,1}}
 \lesssim
 N^{3/2}A
\]
follows by real interpolation between its \(L^2\) and \(L^\infty\)
bounds.  These estimates prove \eqref{eq:refined-speed}.
\end{proof}

\begin{theorem}[Clock--corridor dichotomy]
\label{thm:corridor}
At almost every sufficiently small \(\tau\), at least one of the
following alternatives holds:
\begin{align}
\label{eq:clock-branch}
 \kappa(\tau)
 &\gtrsim
 B_2(\tau)^{-2/3},
 &
 X(\tau)Y(\tau)
 &\gtrsim
 B_2(\tau)^{-4/3};
 \tag{C}\\
\label{eq:corridor-branch}
 \kappa(\tau)H(\tau)
 &\gtrsim1,
 &
 H(\tau)^2X(\tau)Y(\tau)
 &\gtrsim1.
 \tag{R}
\end{align}
\end{theorem}

\begin{proof}
Integrate Lemma~\ref{lem:refined-speed} to the terminal boundary and
put \(N=\kappa(\tau)\):
\[
 \frac{c_0}{2}
 \lesssim
 A\kappa(\tau)^{3/2}B_2(\tau)
 +
 U\kappa(\tau)H(\tau).
\]
At least one term on the right is bounded below by a fixed positive
constant.  If the first pays, then
\(\kappa\gtrsim B_2^{-2/3}\).  If the second pays, then
\(\kappa H\gtrsim1\).  In both cases
Theorem~\ref{thm:cross-motion} supplies the corresponding lower bound
for \(XY\).
\end{proof}

\subsection{Exact q4 record frequencies}

On a retained q4 tail beginning at \(s_j\), the input clock gives
\begin{equation}
\label{eq:record-tails}
 \int_{s_j}^T M\,dt
 \lesssim
 \frac{m_j^{-9/2}}j,
 \qquad
 \int_{s_j}^T M^2\,dt
 \lesssim
 \frac{m_j^{-7/2}}j.
\end{equation}
Put \(\kappa_j:=\kappa(T-s_j)\).  Theorem~\ref{thm:cross-motion} gives
\begin{equation}
\label{eq:mu}
 \boxed{
 \kappa_j
 \gtrsim
 \mu_j:=m_j^{9/4}j^{1/2}.
 }
\end{equation}
The clock branch of Theorem~\ref{thm:corridor} is
\begin{equation}
\label{eq:Lambda}
 \boxed{
 \kappa_j
 \gtrsim
 \Lambda_j:=m_j^{7/3}j^{2/3}.
 }
\end{equation}
Otherwise,
\begin{equation}
\label{eq:record-corridor}
 \boxed{
 \kappa_j
 \int_{s_j}^T M(t)Y(t)\,dt
 \gtrsim1.
 }
\end{equation}
The pre-existing frequency gap is now an exact corridor:
\begin{equation}
\label{eq:frequency-gap}
 \frac{\Lambda_j}{\mu_j}
 \asymp
 m_j^{1/12}j^{1/6}
 \longrightarrow\infty.
\end{equation}
Here \(\mu_j\) is the minimum frequency at which the cross pairing can
enter; \(\Lambda_j\) is where direct high--high nonlinear capacity can
be order one.

\section{Exact scalar endpoint saturation}

Work on \(0<\tau<e^{-e}\), set
\[
 L(\tau):=\log(e/\tau),
\]
and define
\begin{equation}
\label{eq:scalar-functions}
 M_0(\tau)
 :=
 \tau^{-2/11}L(\tau)^{-2/11},
 \qquad
 X_0(\tau)=Y_0(\tau)
 :=
 \tau^{-9/22}L(\tau)^{1/11}.
\end{equation}
Then \(X_0,Y_0\in L^2\), and direct integration gives
\begin{align}
\label{eq:scalar-integrals}
 B_{1,0}(\tau)
 &\asymp
 \tau^{9/11}L(\tau)^{-2/11},
 &
 B_{2,0}(\tau)
 &\asymp
 \tau^{7/11}L(\tau)^{-4/11},
 \\
 H_0(\tau)
 &\asymp
 \tau^{9/22}L(\tau)^{-1/11},
 &
 P_0(\tau)=D_0(\tau)
 &\asymp
 \tau^{2/11}L(\tau)^{2/11}.
\end{align}

Define
\begin{equation}
\label{eq:scalar-kappa}
 \kappa_0(\tau)
 :=
 B_{1,0}(\tau)^{-1/2}
 \asymp
 \tau^{-9/22}L(\tau)^{1/11}.
\end{equation}
The coarse laws and the corridor are saturated:
\begin{equation}
\label{eq:scalar-saturation}
 X_0Y_0
 \asymp
 \kappa_0^2
 \asymp
 B_{1,0}^{-1},
 \qquad
 H_0^2X_0Y_0\asymp1,
\end{equation}
\[
 P_0D_0
 \asymp
 \tau^{4/11}L(\tau)^{4/11}.
\]

The direct clock frequency is
\[
 \Lambda_0(\tau)
 :=
 B_{2,0}(\tau)^{-2/3}
 \asymp
 \tau^{-14/33}L(\tau)^{8/33}.
\]
It strictly exceeds the corridor:
\begin{equation}
\label{eq:below-clock}
 \frac{\kappa_0(\tau)}{\Lambda_0(\tau)}
 \asymp
 \tau^{1/66}L(\tau)^{-5/33}
 \longrightarrow0.
\end{equation}
Indeed,
\begin{equation}
\label{eq:direct-vanishes}
 \kappa_0(\tau)^{3/2}B_{2,0}(\tau)
 \asymp
 \tau^{1/44}L(\tau)^{-5/22}
 \longrightarrow0,
\end{equation}
whereas \(\kappa_0H_0\asymp1\).

Finally,
\begin{equation}
\label{eq:weighted-divergence}
 M_0Y_0^2=\tau^{-1},
 \qquad
 \int_0M_0Y_0^2\,d\tau=\infty.
\end{equation}

\subsection{A synthetic moving front}

Choose a smooth nondecreasing
\(\Phi:[0,\infty)\to[0,1]\) with
\[
 \Phi(z)=0\quad(z\leq1/2),
 \qquad
 \Phi(z)=1\quad(z\geq1),
\]
and define
\begin{equation}
\label{eq:synthetic-front}
 K_N^0(\tau)
 :=
 c_0\Phi\left(\frac{N}{\kappa_0(\tau)}\right).
\end{equation}
Then
\[
 \lim_{\tau\downarrow0}K_N^0(\tau)=0,
 \qquad
 \lim_{N\to\infty}K_N^0(\tau)=c_0,
\]
and its first half-pairing frequency is comparable to \(\kappa_0\).
The derivative is supported where \(N\asymp\kappa_0\), and there
\begin{equation}
\label{eq:synthetic-speed}
 |\partial_\tau K_N^0|
 \lesssim
 \tau^{-1}
 \asymp
 N^2M_0
 \asymp
 NM_0Y_0.
\end{equation}
Thus this front saturates both the coarse speed limit and the corridor
term, while the direct clock term vanishes by
\eqref{eq:direct-vanishes}.

\begin{remark}
This is a scalar temporal ledger.  It has no spatial fields, pressure,
cross measure, or Navier--Stokes dynamics.  It is not a counterexample
to R3C, and \(K_N^0\) is not asserted to be a PDE pairing.  It proves
that the clock, the two energy budgets, all norm-level transition
inequalities derived here, and divergent weighted adjoint dissipation
are mutually consistent at the exact endpoint.
\end{remark}

\section{What changed in this round}

\subsection*{Formal findings, subject to external review}

\begin{enumerate}
\item The terminal cross defect defines a dyadic transition frequency
\(\kappa(\tau)\to\infty\).
\item The exact pressure-free identity gives
\(\kappa^2B_1\gtrsim1\), while its unresolved high-frequency pairing
gives \(XY\gtrsim\kappa^2\).
\item The q4 clock therefore forces the pointwise joint-enstrophy law
\eqref{eq:pointwise-floor}, the terminal tail floor
\eqref{eq:PD-floor}, and the inverse-clock condition
\eqref{eq:inverse-clock}.
\item A sharper Lorentz estimate gives the clock--corridor dichotomy
\eqref{eq:clock-branch}--\eqref{eq:corridor-branch}.
\item The two record frequencies are exactly
\(\mu_j=m_j^{9/4}j^{1/2}\) and
\(\Lambda_j=m_j^{7/3}j^{2/3}\).
\item The scalar ledger \eqref{eq:scalar-functions}, together with the
synthetic front \eqref{eq:synthetic-front}, saturates the corridor, both
speed bounds, and all energy exponents while retaining
\eqref{eq:weighted-divergence}.
\end{enumerate}

\subsection*{Closed shortcut}

The pressure-free cross identity, after taking only the Lorentz norm
bounds and logarithmic clock used here, cannot force the preceding
round's summable amplitude anti-correlation.  Those norm consequences
have the opposite direction: persistent joint enstrophy concentration.
No rearrangement of the proved scalar inequalities can exclude the
exact saturation ledger.  Possible phase or sign information in the
identity has not been discarded as a future route.

\subsection*{Open obligations}

\begin{enumerate}
\item Exclude persistent sub-clock corridor transfer using
same-trajectory phase, sign, genealogy, or projected-pressure structure.
\item Alternatively force the clock branch on enough record blocks and
turn its pointwise \(XY\) floor into a non-reusable residence or
dissipation charge.
\item Prove that successive order-one cross-pairing transitions cannot
recycle one finite primal--adjoint energy reservoir.
\item A finite Lorentz--Zygmund adjoint cost, commutator
equiintegrability, removable total current, strong left \(L^2\)
continuity, or energy equality still closes q4.
\item Treat slower clocks, divergent normalised energy, and the other
Clay alternatives separately.
\end{enumerate}

\begin{conjecture}[No persistent sub-clock cross corridor]
Under the exact q4 hypotheses, one cannot have
\[
 \kappa_j=o(\Lambda_j)
\]
through infinitely many first-record blocks while
\[
 \kappa_j
 \int_{s_j}^TM(t)Y(t)\,dt
 \gtrsim1.
\]
\end{conjecture}

The conjecture is not proved.  Even if it is proved, a residence or
non-reuse theorem may still be needed to exclude the clock branch.  No
Clay alternative is proved.

\end{document}
